% Phase-1 draft abstract, adapted from the thesis proposal abstract.
% TODO: revise, and check Rackham's abstract length limit (550 words).
This dissertation studies how program synthesis and large language
models (LLMs) can rewrite SQL beyond existing rule-based optimizers,
progressing from single-query transformations to full SQL pipeline
refactoring.

SlabCity is a whole-query, synthesis-based SQL rewriter that searches
for semantically equivalent, faster queries without relying on
predefined rules. It defines dataflows for SQL queries and uses them
to stratify the synthesis search space via a novel dataflow-based
scoring function that focuses the search on promising rewrites, within
a counterexample-guided inductive synthesis loop that resorts to SMT
verification only when needed. On curated LeetCode and Calcite
workloads, it rewrites more queries and achieves multi-fold latency
reductions compared to rule-based systems.

GenRewrite targets the same single-query setting but replaces
synthesis-based search with LLM-guided generation anchored by
Natural-Language Rewrite Rules (NLR2s), which capture successful
rewrite strategies in human-readable form and transfer rewriting
knowledge across queries. Across TPC-DS and JOB (including their
SQLStorm-generated variants), GenRewrite attains higher coverage and
larger speedups than both rule-based systems and baseline LLM
prompting.

DAGSmith lifts these ideas from single queries to pipelines of
interdependent SQL models. Given a workload represented as a query
DAG, it combines structural graph analysis, LLM-guided refactoring
with data-backed equivalence checking, learned cost modeling,
materialization tuning via iterative ILP, and frequency-aware
scheduling to produce output-equivalent pipelines that are
significantly cheaper and faster end to end.

Together, these systems trace a path from synthesis to LLMs for SQL
query rewriting beyond rules, delivering substantial performance gains
from individual queries to large, dependency-rich SQL pipelines.
